\subsection{KT-F149-COMP: the two normalizer comparisons agree in scope}
\label{sec:kt-f149-comp}

GRH\_NIGHT AG-C149-NORMAL is an explicit isotrivial instance
of the general characteristic-zero comparison K17
F149-RELATIVE, and it also verifies the homotopy injection
which F149-RELATIVE retains as a hypothesis. Here is the
precise comparison, including the line-bundle convention.

Let $S/k$ be smooth, proper and geometrically connected in
characteristic zero, and choose
$\tau:\mathcal L^{\otimes n}\simeq\mathcal O_S$.
AG-C149-NORMAL defines
\[
 U=\{v\in\mathcal L^\times:
                         \tau(v^{\otimes n})\ne1\}.
\]
The zero-cusp normal line is $\mathcal L$; its conormal
line in F149-RELATIVE's convention is $\mathcal L^\vee$.
Thus their integral Chern classes differ by the stated
normal/conormal sign, not by an arithmetic discrepancy.

To verify the injection, pass to the finite \'etale torsor
$T=\{e:\tau(e^{\otimes n})=1\}$, and, after geometric
base change, choose a connected component. The pulled-back
family is the product of that proper component with
$\mathbf P^1\setminus(\{0,\infty\}\cup\mu_n)$.
SGA~1, Expos\'e X, Corollaire~1.7 injects the fibre
fundamental group into this product group. The finite
\'etale map from the product to $U$ injects its fundamental
group into $\pi_1(U)$. Right exactness is the relative
homotopy sequence of SGA~1, Expos\'e XIII, Proposition~4.1
and Examples~4.4. This proves exactly the sequence
(R149.1) required by F149-RELATIVE in this family.

F149-RELATIVE consequently identifies the zero-cusp
normalizer with $\pi_1(\mathcal L^\times)$.
AG-C149-NORMAL supplies an explicit map in the reverse
direction: the filling map $U\to\mathcal L^\times$
restricts to an isomorphism on the normalizer. The two
descriptions identify the same extension over $\pi_1(S)$,
with the same fibre inertia after the normal/conormal
convention is fixed. Tangential choices and conjugacies
must be transported together; an arbitrary automorphism
of an extension acting identically on kernel and quotient
need not itself be the identity.

AG-C149-FULLLIFT concerns pullback along each arithmetic
section $G_k\to\pi_1(S)$, not a splitting over all of
$\pi_1(S)$. Hence its coexistence with the nonsplit
normalizer in AG-C149-CUSP is consistent with
F149-RELATIVE. For the example
$S:y^2=x^3-x$ over $\mathbf Q$ and its nontrivial
two-torsion line, the global integral class is nonzero,
while its pullback along each arithmetic section vanishes.
Neither statement addresses the independent Ihara kernel
question LM6 A.45.

The complete supplied AG-C149-NORMAL and AG-C149-FULLLIFT
proofs were read, as was F149-RELATIVE. Their source
locators are GRH\_NIGHT.tex, lines 9063--9227, and
K17.tex, lines 13494--13644. The elliptic class computation
AG-C149-ELLIPTIC is carried from the intake, not
independently recomputed here. This comparison concerns
cote~149 scans~79--82; scan~85's geometric-subtorsor
recognition is the separate assertion reviewed above.

